\begin{lemma}[Reduction of the theorem of sums to a quadratic test with a strict global maximum]
Let $E$ be a finite-dimensional real inner-product space, let $N\ge1$ and write
$V=\ell^2\bigl(\bigoplus_{j=1}^NE\bigr)$ for the $\ell^2$-product of $N$ copies
of $E$, whose elements are written $Y=(y_j)_{j=1}^N$ and whose norm satisfies
$|Y|^2=\sum_j|y_j|^2$. For each $j$ let $O_j\subseteq E$ be open and let
$g_j:E\to\R$ be continuous on $O_j$. Let $\Psi:V\to\R$, let $\widehat Z\in V$
with $\widehat z_j\in O_j$ for every $j$, assume that $\Psi$ is $C^2$ near
$\widehat Z$, and assume that $Y\mapsto\sum_jg_j(y_j)-\Psi(Y)$ has a local
maximum at $\widehat Z$. Write $A=D^2\Psi(\widehat Z)$ and
$\widehat p_j=D_j\Psi(\widehat Z)$ for the $j$-th component of
$\nabla\Psi(\widehat Z)$, and let $\tau>0$. Then there are $a>0$, $C\ge0$ and
functions $G_j:E\to\R$ $(1\le j\le N)$ such that
\begin{enumerate}
\item[(G1)] each $G_j$ is continuous on $E$ and $G_j(w)\le C$ for every $w\in E$;
\item[(G2)] $G_j(w)=g_j(\widehat z_j+w)-g_j(\widehat z_j)
  -\langle\widehat p_j,w\rangle-|w|^4$ whenever $|w|\le a$;
\item[(G3)] for every $W\in V$ with $W\ne0$,
  \[
    \sum_jG_j(w_j)-\tfrac12\bigl(A[W]+\tau|W|^2\bigr)\le-\sum_j|w_j|^4 ,
  \]
  where $A[W]=A(W,W)$.
\end{enumerate}
\end{lemma}
\begin{proof}
Write $\widehat P=\nabla\Psi(\widehat Z)$, so that $\widehat P$ has components
$\widehat p_j$ and
\begin{equation}\label{eq:srqt-inner}
  \langle\widehat P,W\rangle=\sum_j\langle\widehat p_j,w_j\rangle\qquad(W\in V),
\end{equation}
by the definition of the $\ell^2$-product inner product (preamble). Put
\[
  \widehat g_j(w)=g_j(\widehat z_j+w)-g_j(\widehat z_j)
   -\langle\widehat p_j,w\rangle ,
\]
so $\widehat g_j(0)=0$. Let $\bar A:V\to V$ be the self-adjoint operator with
$\langle\bar A\xi,\zeta\rangle=A(\xi,\zeta)$ (it exists and is unique because
$A$ is a symmetric continuous bilinear form on the finite-dimensional
inner-product space $V$, and self-adjointness is the symmetry of $A$), and set
$B_\tau=\bar A+\tau\,\mathrm{id}$, again self-adjoint, so that
\begin{equation}\label{eq:srqt-Btau}
  \langle B_\tau W,W\rangle=A[W]+\tau|W|^2\qquad(W\in V).
\end{equation}

\emph{Step 1: a quadratic majorant near $0$.}
By the second-order Taylor expansion of $\Psi$ at $\widehat Z$ (preamble),
$\Psi(\widehat Z+W)=\Psi(\widehat Z)+\langle\widehat P,W\rangle
+\tfrac12A(W,W)+R(W)$ with $|R(W)|/|W|^2\to0$ as $W\to0$. The local maximum at
$\widehat Z$ gives $\rho'>0$ with
\[
  \sum_jg_j(\widehat z_j+w_j)-\Psi(\widehat Z+W)
   \le\sum_jg_j(\widehat z_j)-\Psi(\widehat Z)\qquad(|W|\le\rho').
\]
Using \eqref{eq:srqt-inner}, for $|W|\le\rho'$,
\[
  \sum_j\widehat g_j(w_j)
   =\sum_jg_j(\widehat z_j+w_j)-\sum_jg_j(\widehat z_j)
     -\langle\widehat P,W\rangle
   \le\Psi(\widehat Z+W)-\Psi(\widehat Z)-\langle\widehat P,W\rangle
   =\tfrac12A(W,W)+R(W).
\]
Since $|R(W)|/|W|^2\to0$, choose $\rho_\tau\in(0,\rho']$ with
$|R(W)|\le(\tau/2)|W|^2$ for $|W|\le\rho_\tau$. Because
$\tfrac12A(W,W)+(\tau/2)|W|^2=\tfrac12\langle B_\tau W,W\rangle$ by
\eqref{eq:srqt-Btau}, we obtain
\begin{equation}\label{eq:srqt-majorant}
  \sum_j\widehat g_j(w_j)\le\tfrac12\langle B_\tau W,W\rangle
  \qquad(|W|\le\rho_\tau),
\end{equation}
with equality at $W=0$, both sides vanishing because $\widehat g_j(0)=0$.

\emph{Step 2: localization.}
Each $O_j$ is open and $\widehat z_j\in O_j$, so for each $j$ there is
$r_j>0$ with $\overline B(\widehat z_j,r_j)\subseteq O_j$; as there are finitely
many $j$ and $N\ge1$, the number $\rho''=\min_jr_j>0$ satisfies
$\overline B(\widehat z_j,\rho'')\subseteq O_j$ for every $j$. Put
\[
  a=\frac{\min(\rho_\tau,\rho'')}{1+\sqrt N}>0 .
\]
Then $a\le\rho''$, so $\widehat g_j$ is defined and continuous on
$\overline B(0,a)$ (being built from $g_j$, which is continuous on $O_j\supseteq
\overline B(\widehat z_j,\rho'')$, and from an affine function); and if
$|w_j|\le a$ for every $j$ then $|W|\le a\sqrt N\le\rho_\tau$, so
\eqref{eq:srqt-majorant} applies to $W$. Let
$M_j=\max_{\overline B(0,a)}|\widehat g_j|$, which is finite by the continuity of
$\widehat g_j$ on the compact set $\overline B(0,a)$ (preamble) and satisfies
$M_j\ge|\widehat g_j(0)|=0$; put $C=\max_jM_j\ge0$.

Let $\pi:E\to\overline B(0,a)$ be the radial retraction, $\pi(w)=w$ for
$|w|\le a$ and $\pi(w)=aw/|w|$ for $|w|>a$; it is continuous, and
$|\pi(w)|\le a$ for every $w$. Write $(s)_+=\max(s,0)$, put
$K=\|B_\tau\|N(a+1)+1$ and define
\[
  G_j(w)=\widehat g_j(\pi(w))-K(|w|-a)_+-K(|w|-a)_+^2-|w|^4 .
\]

\emph{(G1).} Each $G_j$ is continuous on $E$, being a composition and
combination of the continuous maps $\pi$, $\widehat g_j|_{\overline B(0,a)}$,
$w\mapsto(|w|-a)_+$ and $w\mapsto|w|^4$. The three subtracted terms are
nonnegative, so
$G_j(w)\le\widehat g_j(\pi(w))\le M_j\le C$ for every $w$, using
$|\pi(w)|\le a$.

\emph{(G2).} If $|w|\le a$ then $\pi(w)=w$ and $(|w|-a)_+=0$, so
$G_j(w)=\widehat g_j(w)-|w|^4$, which is the asserted identity by the definition
of $\widehat g_j$.

\emph{(G3).} Put $\Phi(W)=\sum_jG_j(w_j)-\tfrac12\langle B_\tau W,W\rangle$; by
\eqref{eq:srqt-Btau} the assertion (G3) is exactly
$\Phi(W)\le-\sum_j|w_j|^4$ for $W\ne0$. Let $W\ne0$, put
$m=\max_j(|w_j|-a)_+\ge0$ and $\Pi W=(\pi(w_j))_j$.

If $m=0$ then $|w_j|\le a$ for every $j$, so $\pi(w_j)=w_j$ and all penalty
terms vanish; since $|W|\le a\sqrt N\le\rho_\tau$, \eqref{eq:srqt-majorant}
gives $\sum_j\widehat g_j(w_j)\le\tfrac12\langle B_\tau W,W\rangle$ and hence
$\Phi(W)\le-\sum_j|w_j|^4$.

If $m>0$, then $|\pi(w_j)|\le a$ for every $j$ gives
$|\Pi W|\le a\sqrt N\le\rho_\tau$, so \eqref{eq:srqt-majorant} applied to
$\Pi W$ gives
$\sum_j\widehat g_j(\pi(w_j))\le\tfrac12\langle B_\tau\Pi W,\Pi W\rangle$.
Moreover $\sum_j(|w_j|-a)_+\ge m$ and $\sum_j(|w_j|-a)_+^2\ge m^2$, since every
summand is nonnegative and one of them equals $m$, resp.\ $m^2$. Hence
\[
  \Phi(W)\le\tfrac12\bigl[\langle B_\tau\Pi W,\Pi W\rangle
     -\langle B_\tau W,W\rangle\bigr]-Km-Km^2-\sum_j|w_j|^4 .
\]
Since $B_\tau$ is self-adjoint,
\[
  \langle B_\tau W,W\rangle-\langle B_\tau\Pi W,\Pi W\rangle
   =\langle B_\tau(W-\Pi W),W+\Pi W\rangle .
\]
Here $|W-\Pi W|^2=\sum_j|w_j-\pi(w_j)|^2
=\sum_j\bigl((|w_j|-a)_+\bigr)^2\le Nm^2$, because
$|w_j-\pi(w_j)|=(|w_j|-a)_+$ by the definition of $\pi$ (it is $0$ for
$|w_j|\le a$ and $|w_j|-a$ otherwise); so $|W-\Pi W|\le\sqrt Nm$, and
$|W+\Pi W|\le|W|+|\Pi W|$. From $|w_j|\le a+m$ for every $j$ we get
$|W|\le(a+m)\sqrt N$, and $|\Pi W|\le a\sqrt N$, so
$|W+\Pi W|\le2a\sqrt N+\sqrt Nm$. By the Cauchy--Schwarz inequality and the
definition of the operator norm,
\[
  \tfrac12\bigl|\langle B_\tau\Pi W,\Pi W\rangle-\langle B_\tau W,W\rangle\bigr|
   \le\tfrac12\|B_\tau\|\sqrt Nm\bigl(2a\sqrt N+\sqrt Nm\bigr)
   =\|B_\tau\|Nam+\tfrac12\|B_\tau\|Nm^2,
\]
and therefore
\[
  \Phi(W)\le m\bigl(\|B_\tau\|Na-K\bigr)+m^2\bigl(\tfrac12\|B_\tau\|N-K\bigr)
     -\sum_j|w_j|^4\le-\sum_j|w_j|^4,
\]
because $K=\|B_\tau\|N(a+1)+1\ge\|B_\tau\|Na$ and
$K\ge\|B_\tau\|N\ge\tfrac12\|B_\tau\|N$, while $m>0$ and $m^2>0$. This proves
(G3) in both cases.
\end{proof}
